\chapter{The Death of Spacetime Curvature}
\label{ch:spacetime-curvature}

\begin{quote}
\emph{This chapter proposes a mechanical substrate interpretation that reproduces all observational predictions of General Relativity. We do not claim GR is wrong---its mathematics is extraordinarily successful. We claim the paradox it creates is resolved when the underlying mechanism is supplied.}
\end{quote}

\section{The Paradox}

General Relativity describes gravity as the curvature of a four-dimensional spacetime manifold. Objects follow geodesics through this curved background. The mathematics is extraordinarily successful: GPS satellites carry relativistic corrections accurate to nanoseconds; gravitational lensing predictions match observations to arc-second precision; gravitational wave detections match GR waveforms to extraordinary fidelity.

And yet: \emph{what is curvature?}

GR describes the \emph{effect} of gravity with unmatched precision, but provides no \emph{mechanism}. How does mass ``tell space to curve''? What physically changes when a manifold acquires curvature? The ``fabric of spacetime'' is a metaphor for a mathematical transformation, not a physical entity with demonstrable mechanical properties.

This is not a criticism of GR. It is a categorisation: GR is a phenomenological description, not a mechanical explanation.


\section{The Resolution}

The paradox dissolves when gravity is given a mechanical cause.

\textbf{SDT:} Space is a tangible, pressurised medium (the spation field). Mass displaces this medium, creating pressure gradients. An orbiting body follows a path of minimum resistance through the pressure landscape---the physical realisation of a geodesic.

\begin{center}
\begin{tabular}{ll}
\toprule
\textbf{GR Description} & \textbf{SDT Mechanism} \\
\midrule
``Curvature of spacetime'' & Pressure gradient in spation \\
``Geodesic'' & Path of least resistance \\
``Mass tells space how to curve'' & Mass displaces spation, creating pressure \\
``Space tells mass how to move'' & Pressure gradient accelerates mass \\
\bottomrule
\end{tabular}
\end{center}

The mathematics of curved manifolds is a \emph{correct description} of the pressure geometry. GR is not wrong; it is incomplete. It describes the shape of the pressure field without identifying the existence of the field itself.

\subsection{The Equivalence}

The SDT velocity formula $v = (c/\kop)\sqrt{R/r}$ and the GR orbital velocity $v = \sqrt{GM/r}$ are not competing predictions---they are \textbf{algebraically identical} through the equivalence $GM = c^2 R/\kop^2$.

Every prediction of GR is preserved. Gravitational lensing, frame-dragging, gravitational waves, and perihelion precession all follow from the pressure-gradient model. What changes is the \emph{interpretation}: curvature is not a property of an abstract manifold, but the geometry of a physical pressure field.

\subsection{When the Mechanism Exists, the Paradox Dies}

The ``paradox'' of spacetime curvature---how does matter physically curve a mathematical abstraction?---was never a paradox. It was a symptom of missing information. The moment you identify the mechanical cause (pressure gradients in a tangible medium), the question answers itself, and the paradox ceases to exist.

\textbf{The spacetime curvature paradox is dead.}
